\section*{Chapitre \ref{ch:m1:the-sell-side} --- Le sell side}

\begin{solution}{exo:m1:the-sell-side:1}
$30 \times 10^6 \times 0.015 \times \sqrt{3/3} = \$450\,000$~: trois jours de
vente linéaire portent le risque d'une journée sur la position complète.
\end{solution}

\begin{solution}{exo:m1:the-sell-side:2}
$T = 500\,000/(0.1 \times 5\,000\,000) = 1$ jour. $I = 0.7 \times 0.015 \times
\sqrt{0.1} = 0.33\%$.
\end{solution}

\begin{solution}{exo:m1:the-sell-side:3}
Frais~: $0.035 \times 800 = 28$~millions de dollars. L'émetteur a vendu 20 millions d'actions à
40~\$ que le marché valorisait à 46~\$ ce soir-là~: 120~millions de dollars laissés sur la
table, perçus par les investisseurs auxquels des actions avaient été attribuées --- les
clients des banques, ce qui explique que l'attribution est elle-même un actif de \omterm{def:m1:the-sell-side:flow}{franchise}.
\end{solution}

\begin{solution}{exo:m1:the-sell-side:4}
$T = 3/(0.1\times 6) = 5$ jours. $I = 0.7 \times 0.025 \times \sqrt{0.5} =
1.24\%$~; risque $= 0.025\sqrt{5/3} = 3.23\%$~; $d = 1.24 + 1.645 \times 3.23 =
6.55\%$, soit 1.31~\$ par action.
\end{solution}

\begin{solution}{exo:m1:the-sell-side:5}
$d_z = \sigma\sqrt{Q/V}\,\bigl(\kappa + z/\sqrt{3\pi}\bigr)$, et le crochet
ne dépend pas de $Q$. Pour une taille neuf fois plus grande~: la décote est triplée, soit 6~\%.
\end{solution}

\begin{solution}{exo:m1:the-sell-side:6}
(a) Le financement $4.251 + 7.195 = 11.4$~milliards de dollars, 28~\% de la division~;
l'intermédiation $19.6$~milliards de dollars, 47~\%~; les honoraires $9.3$~milliards de dollars, 23~\% (le
petit reste correspond à d'autres revenus). Le financement représente 37~\% du total FICC et actions.
(b) Le financement relève du carry~; l'intermédiation relève du spread, avec une part de
P\&L de position~; les honoraires de banque d'investissement sont des honoraires.
\end{solution}

\begin{solution}{exo:m1:the-sell-side:7}
(a) La volatilité résiduelle vaut $0.02\sqrt{0.6} = 1.55\%$~; risque
$1.55\% \times \sqrt{2/3} = 1.26\%$~; $d = 0.63 + 1.645 \times 1.26 =
2.71\%$. (b) La couverture abaisse la décote d'équilibre de 61 points de base pour
un coût de 2~: elle vaut trente fois ce qu'elle coûte, ce qui explique que chaque \omterm{def:m1:the-sell-side:block}{transaction de bloc}
soit couverte dans les secondes qui suivent la transaction.
\end{solution}

\begin{solution}{exo:m1:the-sell-side:8}
Dans le modèle~: la règle de la racine carrée était énoncée comme insensible à la durée
\emph{dans des limites raisonnables}~; à la moitié du volume, le dealer est le
marché, ses ventes sont visibles de tous et l'impact réalisé s'élève bien
au-dessus de $\kappa\sigma\sqrt{Q/V}$, si bien que le premier terme n'est plus constant en
$\pi$. Hors du modèle~: le volume n'est pas une donnée. D'autres intervenants détectent
le vendeur et se retirent ou vendent devant lui, de sorte que le volume par rapport auquel
50~\% est mesuré se réduit, et la trajectoire des cours acquiert une dérive défavorable au
dealer que la marche aléatoire sans dérive exclut.
\end{solution}

\begin{solution}{pb:m1:the-sell-side:1}
\textbf{1.} 100~millions de dollars~; 50~\% du volume quotidien.
\textbf{2.} $T = 4/(0.1\times 8) = 5$ jours.
\textbf{3.} $0.7\times 0.018\times\sqrt{0.5} = 0.89\%$, 891\,000~\$.
\textbf{4.} $0.018\sqrt{5/3} = 2.32\%$, 2.32~millions de dollars.
\textbf{5.} $d = 0.89 + 1.645\times 2.32 = 4.71\%$~; cours acheteur 23.82~\$.
\textbf{6.} $T = 2.5$ jours~; $d = 0.89 + 1.645 \times 1.64 = 3.59\%$.
\textbf{7.} Volatilité résiduelle $0.018/\sqrt2 = 1.27\%$~; risque $1.16\%$~;
$d = 0.89 + 1.645\times 1.16 = 2.80\%$.
\textbf{8.} La décote moins l'impact~: $1.91\%$ de 100~millions de dollars,
1.91~millions de dollars.
\textbf{9.} $0.89 + 1.16 = 2.05\%$~; elle perd de l'argent avec une probabilité
$1-\Phi(1) = 16\%$.
\textbf{10.} Coût espéré $0.89 + 0.04 = 0.93\%$~; écart-type
$0.018\sqrt{2.5/3} = 1.64\%$ (le fonds ne se couvre pas).
\textbf{11.} $2.80 - 0.93 = 1.87\%$, 1.87~millions de dollars~: le fonds achète la
suppression d'un écart-type de 1.64~millions de dollars et de tout l'effort d'exécution.
\textbf{12.} Son indice de référence est le cours de décision~; face à lui, une exécution
d'agence peut dériver de 3~\% ou 4~\% un mauvais jour, dont elle répond
personnellement, alors que l'écart d'exécution du risk bid est connu, convenu et définitif à
11:41.
\textbf{13.} Ajouter 0.4~\% (les \omterm{def:m1:the-sell-side:block}{transactions de bloc} couvertes perdent exactement le mouvement par rapport à
l'indice)~: c'est le coût de \omterm{def:m1:what-a-trading-firm-does:adverse-selection}{sélection adverse} du flux du client.
\textbf{14.} $150\times10^6 \times 0.12 \times 0.10 \times 3/252 = \$21\,429$.
\textbf{15.} Non~: le capital coûte environ 1~\% du profit espéré. Le risque, et non le
capital, fixe le prix d'une position de trois jours~; le capital contraint les positions détenues pendant
des mois.
\textbf{16.} Coter un cours acheteur au coût d'impact, soit une décote de 0.89~\%, renonce à la
totalité du profit espéré de 1.9~millions de dollars, tout en conservant une perte espérée nulle~;
en deçà, le desk finance la relation sur les 3~millions de dollars, et
ne devrait le faire que si perdre le bloc revient à perdre le client.
\textbf{17.} $-3$~millions de dollars sur l'action, $+\$1$~million sur les contrats à terme~:
$-2$~millions de dollars, soit déjà 1.7 écarts-types du dénouement couvert.
\textbf{18.} Que la position a été prise pour répondre à la demande d'un client et que
les limites d'inventaire et de risque du desk sont conçues pour ne pas dépasser les
demandes raisonnablement attendues à court terme de ses clients~: une demande documentée
d'un client, des limites calibrées sur les flux passés, et un dénouement cohérent avec elles.
\textbf{19.} $2.80 + 0.40 = 3.20\%$, soit \textbf{80 cents par action}.
\textbf{20.} Le client peut vendre au cours acheteur ou se retirer, et il vend plus
volontiers lorsqu'il croit que le cours est sur le point de baisser~: le dealer a vendu
un put dont l'exercice est corrélé aux mauvaises nouvelles.
\end{solution}

\begin{solution}{iq:m1:the-sell-side:1}
Le \omterm{def:m1:the-sell-side:sales-trader}{sales-trader} traite l'ordre du client sur le marché en agent~: le client
détient chaque exécution et le risque d'exécution, et le \omterm{def:m1:the-sell-side:sales-trader}{sales-trader} est
jugé sur la qualité d'exécution par rapport à un indice de référence. Le trader achète ou vend
contre le client en principal~: une fois le prix convenu, le P\&L de la
position revient au book du trader.
\iqlookfor{principal contre agent, formulé en termes de qui détient le risque après
la fin de l'appel téléphonique.}
\end{solution}

\begin{solution}{iq:m1:the-sell-side:2}
Une position qui se réduit linéairement a la même variance qu'une position constante détenue pendant
un tiers de cette durée~: son écart-type vaut $\sqrt{1/3} \approx 58\%$ de celui
d'une détention totale pendant neuf jours --- soit le risque de la détenir entièrement pendant trois jours.
\iqlookfor{la règle $T/3$ de mémoire ou dérivée en une minute.}
\end{solution}

\begin{solution}{iq:m1:the-sell-side:3}
Quelle est sa taille par rapport au volume quotidien, et quelle est la volatilité (la
décote mécanique)~; puis-je le couvrir, et avec quoi~; pourquoi ce client
vend, et ce qui s'est passé après ses blocs précédents~; y a-t-il une
événement (résultats, changement d'indice) pendant ma fenêtre de dénouement~; quelle est ma position
actuelle et ma limite sur le titre et sur le secteur~; et ce que vaut le client
pour la banque.
\iqlookfor{la \omterm{def:m1:what-a-trading-firm-does:adverse-selection}{sélection adverse} et l'inventaire existant mentionnés, et pas seulement la
formule.}
\end{solution}

\begin{solution}{iq:m1:the-sell-side:4}
Le trading pour compte propre a été interdit aux banques et les exigences de
fonds propres liées au risque de marché ont augmenté, si bien que le risque directionnel lent rapporte trop peu par unité de
capital. Le financement, piloté par la clientèle et collatéralisé, consomme peu
de fonds propres au titre du risque de marché par dollar de revenu, se reproduit chaque jour et a grandi avec
le secteur des \omterm{def:m1:what-a-trading-firm-does:hedge-fund}{hedge funds} que d'anciens traders de banques sont allés bâtir.
\iqlookfor{le rendement sur fonds propres comme idée directrice, ainsi que la
cause réglementaire.}
\end{solution}

\begin{solution}{iq:m1:the-sell-side:5}
Un bloc de 4 millions d'actions coûte $2c$ par action si 1 million coûte $c$~: un total
de $8c$ millions. Quatre blocs de 1 million coûtent $4c$ millions~: soit moitié moins
--- \emph{si} l'impact est totalement oublié d'un jour à l'autre et que le
marché n'en déduit pas que trois blocs supplémentaires suivent. Si l'impact est permanent
ou si le programme est détecté, les blocs suivants sont fixés à partir d'un cours
plus bas et l'avantage se réduit vers zéro. Le client porte aussi le risque de prix
des actions invendues pendant trois jours supplémentaires.
\iqlookfor{l'argument de concavité, puis les deux réserves (décroissance de l'impact,
fuite d'information) sans qu'on le demande.}
\end{solution}

\begin{solution}{iq:m1:the-sell-side:6}
$\Var\bigl(\tfrac1T\int_0^T W_t\dd t\bigr) = \tfrac{1}{T^2}\iint
\min(s,t)\dd s\dd t = \tfrac{2}{T^2}\int_0^T\!\int_0^t s\dd s\dd t =
\tfrac{2}{T^2}\cdot\tfrac{T^3}{6} = T/3$ (fois $\sigma^2$).
\iqlookfor{la covariance $\min(s,t)$ et une intégrale double propre~; bonus pour
la méthode d'intégration par parties, $\int_0^T (T-t)\dd W_t$.}
\end{solution}
